%=======================================================================
\section{Modelling}
\label{sec:model}
%=======================================================================

The system comprises three subsystems connected in series by a closed loop of
pressurised water (the secondary fluid): the HTHP, which operates during the
charging period of duration $t_\mathrm{ch}$; the ORC, which operates during
the discharging period of duration $t_\mathrm{dc}$; and a field of $\Nw$
identical borehole batteries, through which the water flows in one direction
while charging and in the opposite direction while discharging. The two
thermodynamic cycles are modelled as steady-flow cycles evaluated once per
design point (Section~\ref{sec:cycles}); the borehole batteries are modelled
as a transient, one-dimensional heat exchanger marched through repeated
charge--discharge cycles (Sections~\ref{sec:BB}--\ref{sec:pumping}). The two
models exchange the temperatures and mass flow rates of the secondary fluid
and a small number of scalar efficiencies.

%-----------------------------------------------------------------------
\subsection{P2H2P thermodynamic cycles}
\label{sec:cycles}
%-----------------------------------------------------------------------

The P2H2P cycles consist of a two-stage ORC with vapour reheating and
regeneration and a two-stage HTHP with a flash-gas separator and two
internal heat exchangers (regenerators), Fig.~\ref{fig:layout}. Both operate
with cyclopentane, a dry fluid whose saturated-vapour line has a positive
slope on the $T$--$s$ diagram. State points are labelled with a numeral and a
suffix identifying the stream: \textsf{e} for the ORC working fluid,
\textsf{h} for the HTHP refrigerant, \textsf{c} and \textsf{d} for the
secondary water during charging and discharging, and \textsf{a} for the
heat-pump source stream.

\begin{figure}[htbp]
  \centering
  \begin{subfigure}{0.85\linewidth}\centering
    \figorbox{\linewidth}{Block_diag_ch.png}
    \subcaption{Charging period}
  \end{subfigure}\\[4pt]
  \begin{subfigure}{0.85\linewidth}\centering
    \figorbox{\linewidth}{Block_diag_dc.png}
    \subcaption{Discharging period}
  \end{subfigure}
  \caption{Components and state points of the two-stage P2H2P energy
  conversion system. \chk{the IHTC-18 diagrams show $T_\mathrm{2d}$ and
  $T_\mathrm{3d}$ with the old discharge closure; relabel the borehole inlet
  and outlet if needed, and add the ORC evaporator composite curves as an
  inset or separate figure.}}
  \label{fig:layout}
\end{figure}

The following idealisations apply to both cycles. Heat exchangers and
connecting lines are free of pressure drop and heat loss, throttling is
isenthalpic, and the thermophysical properties of all fluids are obtained from
CoolProp \cite{CoolProp}. Compression, pumping and expansion are adiabatic but
irreversible, and are characterised by isentropic efficiencies applied to the
state points themselves: for an expansion from state $a$ to pressure $p_b$
and for a compression (or pumping) from state $a$ to pressure $p_b$,
respectively,
\begin{equation}
  h_b=h_a-\eta_{s,t}\left(h_a-h_{b,s}\right),
  \qquad
  h_b=h_a+\frac{h_{b,s}-h_a}{\eta_{s}},
  \qquad
  h_{b,s}=h\left(p_b,s_a\right),
  \label{eq:eta-s}
\end{equation}
where $h_{b,s}$ is the enthalpy at the end of the isentropic process and
$\eta_{s}$ stands for $\eta_{s,p}$ (ORC pumps) or $\eta_{s,c}$ (HTHP
compressors). $\etaORC$ and $\COP$ are therefore the efficiencies of the
irreversible cycles, and only the electrical conversion efficiencies of the
generator and the motor are applied outside the working fluid
(Section~\ref{sec:budget}). Applying the isentropic efficiencies to the state
points, rather than as multipliers on the work of an internally reversible
cycle, has two consequences. First, the states downstream of each machine
move: the turbine exhausts are hotter, which changes the reheat duty, the
feed-heater split and the evaporator composite curve, and the HTHP suction
states change as described in Section~\ref{sec:hthp}. Second, the entropy
generated in each machine appears in the state points, which is a
prerequisite for a component-wise exergy balance; an isentropic machine
destroys no exergy by construction. \chk{v0.12 reference case: the
multiplier treatment underestimates $\COP$ by \SI{5.4}{\percent} (2.346
against 2.473) and $\eta_\mathrm{ORC}\eta_\mathrm{turb}$ by
\SI{0.2}{\percent}; state this if the comparison with Ref.~\cite{KoIHTC2026}
is discussed.}

\subsubsection{Organic Rankine cycle}
\label{sec:orc}

The ORC (Fig.~\ref{fig:layout}b) combines vapour reheating, which raises the
mean temperature of heat addition and the specific work, with a closed
regenerative feed heater supplied by a bleed from the high-pressure turbine,
which reduces the external heat required in the evaporator. Its state points
are defined as follows. The condensing temperature $T_\mathrm{1e}$ and the
evaporating temperature $T_\mathrm{3e}$ fix the two pressure levels,
\begin{equation}
  p_\mathrm{cond,e}=p_\mathrm{sat}(T_\mathrm{1e}),\qquad
  p_\mathrm{evap,e}=p_\mathrm{sat}(T_\mathrm{3e}),\qquad
  p_\mathrm{int,e}=\tfrac12\left(p_\mathrm{evap,e}+p_\mathrm{cond,e}\right),
  \label{eq:orc-p}
\end{equation}
where the intermediate (reheat and bleed) pressure $p_\mathrm{int,e}$ is taken
as the arithmetic mean of the two. State~3e is saturated vapour at
$p_\mathrm{evap,e}$ (high-pressure turbine inlet); 5e is the outlet of the
expansion of 3e to $p_\mathrm{int,e}$ with efficiency $\eta_{s,t}$,
Eq.~\eqref{eq:eta-s}; 6e is the reheated vapour at $p_\mathrm{int,e}$ and
$T_\mathrm{6e}=T_\mathrm{3e}$ (reheat to the live-steam temperature); 4e is the
outlet of the expansion of 6e to $p_\mathrm{cond,e}$ with efficiency
$\eta_{s,t}$; 1e is saturated liquid at $p_\mathrm{cond,e}$; 2e is the outlet
of pump~1, compressing 1e to $p_\mathrm{evap,e}$ with efficiency
$\eta_{s,p}$; 7e is saturated liquid at $p_\mathrm{int,e}$ (drain of the feed
heater); 8e is the outlet of pump~2, compressing 7e to $p_\mathrm{evap,e}$
with efficiency $\eta_{s,p}$; and 9e is the feed liquid leaving the heater at
$p_\mathrm{evap,e}$ and $T_\mathrm{9e}=T_\mathrm{7e}$, i.e.\ with zero
terminal approach \chk{the feed-heater duty is therefore an upper bound}.

Per unit mass flowing through the evaporator, a fraction $1-y$ is bled at
state~5e and condensed in the feed heater, and the complementary fraction $y$
is reheated, expanded to the condenser and returned by pump~1. An energy
balance on the feed heater, $y\,(h_\mathrm{9e}-h_\mathrm{2e})
=(1-y)(h_\mathrm{5e}-h_\mathrm{7e})$, gives
\begin{equation}
  y=\frac{h_\mathrm{5e}-h_\mathrm{7e}}
         {\left(h_\mathrm{9e}-h_\mathrm{2e}\right)
          +\left(h_\mathrm{5e}-h_\mathrm{7e}\right)},
  \label{eq:orc-y}
\end{equation}
and adiabatic mixing of the two liquid streams gives the evaporator inlet,
\begin{equation}
  h_\mathrm{10e}=y\,h_\mathrm{9e}+(1-y)\,h_\mathrm{8e}.
  \label{eq:orc-h10}
\end{equation}
The heat input per unit evaporator flow is supplied in the main evaporator
($\mathrm{10e}\to\mathrm{3e}$) and in the reheater
($\mathrm{5e}\to\mathrm{6e}$, fraction $y$), and the heat rejected is that of
fraction $y$ in the condenser ($\mathrm{4e}\to\mathrm{1e}$). By an energy
balance on the whole cycle, the net specific work is the difference between the
two, and the thermal efficiency of the cycle is
\begin{equation}
  \etaORC=\frac{\left(h_\mathrm{3e}-h_\mathrm{10e}\right)
     +y\left(h_\mathrm{6e}-h_\mathrm{5e}\right)
     -y\left(h_\mathrm{4e}-h_\mathrm{1e}\right)}
    {\left(h_\mathrm{3e}-h_\mathrm{10e}\right)
     +y\left(h_\mathrm{6e}-h_\mathrm{5e}\right)},
  \label{eq:eta-orc}
\end{equation}
in which the turbine and pump irreversibilities enter through
$h_\mathrm{4e}$, $h_\mathrm{5e}$ (and hence $y$), $h_\mathrm{2e}$ and
$h_\mathrm{8e}$, and the pump work is accounted for implicitly.

\paragraph{Condensing temperature} The condenser rejects heat to a sink at
$T_\mathrm{sink}$, which may rise by $\Delta T_\mathrm{sink,gl}$ across the
condenser; the condensing temperature is set by an approach to the sink
outlet,
\begin{equation}
  T_\mathrm{1e}=T_\mathrm{sink}+\Delta T_\mathrm{sink,gl}+\Delta T_\mathrm{E,sink}.
  \label{eq:T1e}
\end{equation}
With $\Delta T_\mathrm{sink,gl}=0$ the sink is an infinite reservoir.

\paragraph{Evaporating temperature: a pinch constraint}
The evaporator transfers heat from the secondary water, which cools from the
borehole outlet temperature $T_\mathrm{2d}$ to the borehole inlet temperature
$T_\mathrm{3d}$, to two cold streams in parallel: the main stream
($\mathrm{10e}\to\mathrm{3e}$ at $p_\mathrm{evap,e}$: liquid preheating
followed by isothermal boiling at $T_\mathrm{3e}$) and the reheat stream
($\mathrm{5e}\to\mathrm{6e}$ at $p_\mathrm{int,e}$, superheated vapour). The
cold composite curve $T_\mathrm{wf}(Q)$ is constructed from the enthalpies of
both streams as a function of the cumulative duty $Q$ measured from the cold
end, and the hot composite $T_\mathrm{w}(Q)$ from the enthalpy of water at the
loop pressure $p_\mathrm{w}$, scaled to the same total duty. Both are built on
real enthalpy: the boiling plateau is a horizontal segment in $(Q,T)$, and
the specific heat of water varies appreciably over a glide of tens of
kelvin. The minimum temperature difference in the exchanger is
\begin{equation}
  \Delta T_\mathrm{pinch}\left(T_\mathrm{3e}\right)=
  \min_{0\le Q\le Q_\mathrm{tot}}
  \left[T_\mathrm{w}(Q)-T_\mathrm{wf}(Q;T_\mathrm{3e})\right],
  \label{eq:pinch}
\end{equation}
and the evaporating temperature is the largest value that satisfies both a
minimum hot-end approach and a minimum pinch,
\begin{equation}
  T_\mathrm{3e}=\min\Bigl(T_\mathrm{2d}-\Delta T_\mathrm{2D,3E},\;
  \max\bigl\{T:\Delta T_\mathrm{pinch}(T)\ge\Delta T_\mathrm{pinch,ORC}\bigr\}
  \Bigr).
  \label{eq:T3e}
\end{equation}
The hot-end condition alone, which is the natural one to write and the one
used in Ref.~\cite{KoIHTC2026}, is sufficient only when the two streams have
comparable heat-capacity curves. A subcritical cycle with a dry fluid absorbs
the larger part of its heat isothermally, so the pinch lies in the interior of
the exchanger, at the onset of boiling. If a fraction $\phi$ of the cycle heat
input is absorbed at $T_\mathrm{3e}$, and the water specific heat is taken as
constant, the water can supply that heat only from the part of its glide lying
above $T_\mathrm{3e}$, which gives the necessary condition
\begin{equation}
  \frac{T_\mathrm{2d}-T_\mathrm{3e}}{T_\mathrm{2d}-T_\mathrm{3d}}\;\ge\;\phi.
  \label{eq:pinch-necessary}
\end{equation}
Equation~\eqref{eq:pinch-necessary} shows that the evaporating temperature,
and with it $\etaORC$, falls approximately in proportion to the width of the
discharging glide. This couples a storage-side design variable to the power
cycle and is the origin of the interior optimum in the glide discussed in
Section~\ref{sec:results}. The maximisation in Eq.~\eqref{eq:T3e} is performed
by bisection on $T_\mathrm{3e}$ between $T_\mathrm{1e}+\SI{10}{\kelvin}$ and
the hot-end limit.

\subsubsection{High-temperature heat pump}
\label{sec:hthp}

The HTHP (Fig.~\ref{fig:layout}a) is a two-stage vapour-compression cycle with
a flash-gas separator at an intermediate pressure and two regenerators, R1 and
R2. The condensing and evaporating temperatures fix
\begin{equation}
  p_\mathrm{cond,h}=p_\mathrm{sat}(T_\mathrm{2h}),\qquad
  p_\mathrm{evap,h}=p_\mathrm{sat}(T_\mathrm{13h}),\qquad
  p_\mathrm{int,h}=\sqrt{p_\mathrm{evap,h}\,p_\mathrm{cond,h}},
  \label{eq:hp-p}
\end{equation}
with the intermediate pressure taken as the geometric mean. State~2h, the
high-stage compressor outlet, is saturated vapour at $p_\mathrm{cond,h}$.
State~3h, the condenser outlet, is subcooled liquid at $p_\mathrm{cond,h}$
and $T_\mathrm{3h}=T_\mathrm{2h}-\Delta T_\mathrm{sub}$. State~13h, the
evaporator outlet, is saturated vapour at $p_\mathrm{evap,h}$. States~9h and
10h are the saturated liquid and saturated vapour leaving the separator at
$p_\mathrm{int,h}$.

Because cyclopentane is a dry fluid, compression from saturated vapour can
terminate inside the two-phase region. The design therefore fixes the
discharge of each stage and determines its suction state backwards, so that the
high-stage compressor discharges exactly on the dew line at $T_\mathrm{2h}$ and
the condenser has no desuperheating duty. The high-stage suction state
$\mathrm{6h}\equiv\mathrm{11h}$ at $p_\mathrm{int,h}$ is the state whose
compression to $p_\mathrm{cond,h}$ with efficiency $\eta_{s,c}$ ends at
$h_\mathrm{2h}$, and the low-stage suction state~1h at $p_\mathrm{evap,h}$ is
the state whose compression to $p_\mathrm{int,h}$ ends at $h_\mathrm{11h}$, so
that the low-stage discharge coincides with state~11h ($\mathrm{5h}\equiv
\mathrm{11h}$):
\begin{equation}
  h_\mathrm{11h}+\frac{h\left(p_\mathrm{cond,h},s_\mathrm{11h}\right)-h_\mathrm{11h}}{\eta_{s,c}}
  =h_\mathrm{2h},
  \qquad
  h_\mathrm{1h}+\frac{h\left(p_\mathrm{int,h},s_\mathrm{1h}\right)-h_\mathrm{1h}}{\eta_{s,c}}
  =h_\mathrm{11h},
  \label{eq:hp-suction}
\end{equation}
with $s=s(p,h)$ at the respective suction pressure. For $\eta_{s,c}=1$ the
relations reduce to $s_\mathrm{11h}=s_\mathrm{1h}=s_\mathrm{2h}$ and are
explicit; for $\eta_{s,c}<1$ they are implicit in the suction enthalpy and are
solved by bisection, the high stage first. An irreversible compression
ending at the same discharge state requires a colder suction, so the suction
superheat decreases as $\eta_{s,c}$ decreases \chk{\SI{14.3}{\kelvin} at
$\eta_{s,c}=0.85$ against \SI{26.5}{\kelvin} for isentropic compression, v0.12
reference case}. The superheat at 1h and 11h is supplied by the two
regenerators, which cool the condensate on its way to the expansion valves:

\begin{itemize}[leftmargin=1.8em,itemsep=1pt]
\item R1 cools the full condensate flow from 3h to 12h and heats the flashed
  vapour, of mass fraction $x_\mathrm{8h}$, from 10h to 11h;
\item R2 cools the condensate further from 12h to 7h and heats the suction
  flow, of mass fraction $1-x_\mathrm{8h}$, from 13h to 1h.
\end{itemize}

\noindent
The condensate is then throttled from 7h into the separator
($h_\mathrm{8h}=h_\mathrm{7h}$), and the separator liquid is throttled from 9h
into the evaporator ($h_\mathrm{4h}=h_\mathrm{9h}$). Per unit mass of
refrigerant through the condenser, the energy balances on R1, R2 and the
separator are
\begin{align}
  h_\mathrm{12h}&=h_\mathrm{3h}-x_\mathrm{8h}\left(h_\mathrm{11h}-h_\mathrm{10h}\right),
  \label{eq:R1}\\
  h_\mathrm{7h}&=h_\mathrm{12h}-\left(1-x_\mathrm{8h}\right)
                 \left(h_\mathrm{1h}-h_\mathrm{13h}\right),
  \label{eq:R2}\\
  x_\mathrm{8h}&=\frac{h_\mathrm{7h}-h_\mathrm{9h}}{h_\mathrm{10h}-h_\mathrm{9h}}.
  \label{eq:flash}
\end{align}
Equations~\eqref{eq:R1}--\eqref{eq:flash} are linear in $x_\mathrm{8h}$ once
the pressure levels are fixed and are solved by fixed-point iteration, which
contracts strongly and converges to $10^{-12}$ in a few iterations. It is
emphasised that the regenerators are closed by energy rather than by
prescribed effectivenesses: the effectiveness each regenerator requires,
\begin{equation}
  \varepsilon_\mathrm{R1}=\frac{T_\mathrm{11h}-T_\mathrm{10h}}{T_\mathrm{3h}-T_\mathrm{10h}},
  \qquad
  \varepsilon_\mathrm{R2}=\frac{T_\mathrm{1h}-T_\mathrm{13h}}{T_\mathrm{12h}-T_\mathrm{13h}},
  \label{eq:eps-R}
\end{equation}
is an output of the cycle and a feasibility check (a value above unity would
indicate a regenerator that cannot exist). \chk{Earlier versions used the
evaporator-inlet quality $x_\mathrm{4h}$ in place of the flow split
$1-x_\mathrm{8h}$ in Eq.~\eqref{eq:R2}; this violated the energy balance by
\SI{4.8}{\percent} and inflated the COP. State in the paper that the IHTC-18
COP is superseded, or leave implicit.}

The low-stage compressor handles the fraction $1-x_\mathrm{8h}$ from 1h to
5h, and the high-stage compressor the full flow from 6h to 2h, with
$h_\mathrm{6h}=h_\mathrm{5h}$. The heating coefficient of performance of the
cycle is therefore
\begin{equation}
  \COP=\frac{h_\mathrm{2h}-h_\mathrm{3h}}
  {\left(h_\mathrm{2h}-h_\mathrm{1h}\right)
   -x_\mathrm{8h}\left(h_\mathrm{5h}-h_\mathrm{1h}\right)}.
  \label{eq:cop}
\end{equation}

\paragraph{Evaporating and condensing temperatures}
The heat pump draws heat from a source stream (e.g.\ co-produced water) that
enters at $T_\mathrm{3a}=T_\mathrm{source}$ and is cooled by
$\Delta T_\mathrm{3A,4A}$. The evaporating temperature is set by a cold-end
approach to the cooled source, and the condensing temperature by a hot-end
approach to the water leaving the condenser:
\begin{equation}
  T_\mathrm{13h}=T_\mathrm{source}-\Delta T_\mathrm{3A,4A}-\Delta T_\mathrm{pinch,HPE},
  \qquad
  T_\mathrm{2h}=T_\mathrm{3c}+\Delta T_\mathrm{2H,3C}.
  \label{eq:hp-T}
\end{equation}

\subsubsection{Secondary-fluid loop and temperature closure}
\label{sec:closure-T}

The secondary fluid is liquid water at $p_\mathrm{w}=\SI{1}{\mega\pascal}$,
which keeps it subcooled over the full operating range. During charging it is
heated in the HTHP condenser from $T_\mathrm{2c}$ to $T_\mathrm{3c}$ and
enters the borehole field directly, so that the borehole charging inlet is
$T_\mathrm{4c}=T_\mathrm{3c}$ and the borehole charging outlet is
$T_\mathrm{2c}$. During discharging it enters the field at $T_\mathrm{3d}$,
leaves at $T_\mathrm{2d}$, and is cooled in the ORC evaporator back to
$T_\mathrm{3d}$.

Of the four exchanger terminal temperatures of the borehole field, only the
two \emph{inlets} are boundary conditions of the borehole model; the two
outlets are results of the heat transfer and cannot be prescribed without
over-determining the problem. The inlets are set by approaches to the
melting temperatures of the first PCM layer each half-cycle meets,
\begin{equation}
  T_\mathrm{4c}=T_{m}^\mathrm{top}+\Delta T_\mathrm{4C,M},
  \qquad
  T_\mathrm{3d}=T_{m}^\mathrm{bot}-\Delta T_\mathrm{M,1D},
  \label{eq:inlets}
\end{equation}
where $T_{m}^\mathrm{top}$ and $T_{m}^\mathrm{bot}$ are the highest and lowest
melting temperatures of the cascade (Section~\ref{sec:cascade}); the charging
superheat $\Delta T_\mathrm{4C,M}$ drives melting and the discharging
subcooling $\Delta T_\mathrm{M,1D}$ drives freezing.

Three temperature spans of similar magnitude appear in the loop, and they are
not equally free. The charging glide $\Delta T_\mathrm{3C,2C}
=T_\mathrm{3c}-T_\mathrm{2c}$ across the HTHP condenser is specified, and
since nothing lies between the condenser and the borehole the water falls by
the same amount across the field during charging; the discharging glide
$\Delta T_\mathrm{3D,2D}=T_\mathrm{2d}-T_\mathrm{3d}$ across the ORC evaporator
is specified independently and is likewise equal to the rise across the field
during discharging. The cascade span (Section~\ref{sec:cascade}) is a third,
independent specification. The first two fix the two mass flow rates
(Section~\ref{sec:budget}); the outlet temperatures that the field actually
produces are reported and compared with them (Section~\ref{sec:css}).

Because $T_\mathrm{3e}$ depends on $T_\mathrm{2d}$ through
Eq.~\eqref{eq:T3e}, and $T_\mathrm{2d}$ is a result of the borehole model, the
ORC is first evaluated at the provisional outlet
$T_\mathrm{2d}^{(0)}=T_\mathrm{3d}+\Delta T_\mathrm{3D,2D}$ and then once more
at the flow-averaged outlet that the store produces
(Section~\ref{sec:procedure}). Table~\ref{tab:closure} collects the closure
relations.

\begin{table}[htbp]
\centering
\caption{Temperature closure of the ORC, HTHP and secondary-fluid loop.}
\label{tab:closure}
\small
\resizebox{\linewidth}{!}{%
\begin{tabular}{@{}l l l@{}}
\toprule
parameter & relation & role \\
\midrule
$\Delta T_\mathrm{E,sink}$, $\Delta T_\mathrm{sink,gl}$ &
  $T_\mathrm{1e}=T_\mathrm{sink}+\Delta T_\mathrm{sink,gl}+\Delta T_\mathrm{E,sink}$ &
  ORC condensing temperature\\
$\Delta T_\mathrm{pinch,ORC}$, $\Delta T_\mathrm{2D,3E}$ &
  Eq.~\eqref{eq:T3e} & ORC evaporating temperature\\
$\Delta T_\mathrm{M,1D}$ &
  $T_\mathrm{3d}=T_{m}^\mathrm{bot}-\Delta T_\mathrm{M,1D}$ &
  borehole discharging inlet\\
$\Delta T_\mathrm{3D,2D}$ &
  $T_\mathrm{2d}^{(0)}=T_\mathrm{3d}+\Delta T_\mathrm{3D,2D}$ &
  discharging flow; provisional outlet\\
$\Delta T_\mathrm{4C,M}$ &
  $T_\mathrm{4c}=T_\mathrm{3c}=T_{m}^\mathrm{top}+\Delta T_\mathrm{4C,M}$ &
  borehole charging inlet\\
$\Delta T_\mathrm{3C,2C}$ &
  $T_\mathrm{2c}=T_\mathrm{3c}-\Delta T_\mathrm{3C,2C}$ &
  charging flow; nominal outlet\\
$\Delta T_\mathrm{2H,3C}$ &
  $T_\mathrm{2h}=T_\mathrm{3c}+\Delta T_\mathrm{2H,3C}$ &
  HTHP condensing temperature\\
$\Delta T_\mathrm{sub}$ &
  $T_\mathrm{3h}=T_\mathrm{2h}-\Delta T_\mathrm{sub}$ &
  condensate subcooling\\
$\Delta T_\mathrm{3A,4A}$, $\Delta T_\mathrm{pinch,HPE}$ &
  $T_\mathrm{13h}=T_\mathrm{source}-\Delta T_\mathrm{3A,4A}-\Delta T_\mathrm{pinch,HPE}$ &
  HTHP evaporating temperature\\
\bottomrule
\end{tabular}}
\end{table}

\subsubsection{Energy budget and performance indicators}
\label{sec:budget}

For a specified net electrical output $\dot W_\mathrm{el,out}$ over
$t_\mathrm{dc}$, the heat rate and energy the ORC requires from the field are
\begin{equation}
  \dot Q_\mathrm{in,ORC}=\frac{\dot W_\mathrm{el,out}}
     {\etaORC\,\eta_\mathrm{gen}},
  \qquad
  \Delta E_\mathrm{in,ORC}=\dot Q_\mathrm{in,ORC}\,t_\mathrm{dc},
  \label{eq:E-orc}
\end{equation}
where $\eta_\mathrm{gen}$ is the generator efficiency. The isentropic
efficiencies of the machines are already contained in $\etaORC$ and $\COP$
through Eq.~\eqref{eq:eta-s}, and applying them again here would count them
twice. The energy
the HTHP must deliver to the field over $t_\mathrm{ch}$, and the
corresponding heat rate, are
\begin{equation}
  \Delta E_\mathrm{out,HP}=\left(1+\lambda\right)\Delta E_\mathrm{in,ORC},
  \qquad
  \dot Q_\mathrm{out,HP}=\frac{\Delta E_\mathrm{out,HP}}{t_\mathrm{ch}},
  \label{eq:E-hp}
\end{equation}
and the electrical power drawn by the compressors is
\begin{equation}
  \dot W_\mathrm{el,in}=\frac{\dot Q_\mathrm{out,HP}}
    {\COP\,\eta_\mathrm{mot}},
  \label{eq:W-in}
\end{equation}
with $\eta_\mathrm{mot}$ the efficiency of the compressor motors. The storage-loss allowance $\lambda$
was set to \SI{5}{\percent} in Ref.~\cite{KoIHTC2026}; in the present model
the store is adiabatic and at cyclic steady state it returns exactly the
energy it absorbs, so $\lambda=0$ (Section~\ref{sec:css}).

The mass flow rates of water per hairpin follow from the prescribed glides,
\begin{gather}
  \md_\mathrm{ch}=\frac{\dot Q_\mathrm{out,HP}}
     {\left[h_\mathrm{w}(T_\mathrm{3c})-h_\mathrm{w}(T_\mathrm{2c})\right]\Nw\,n_{t}},
  \notag\\
  \md_\mathrm{dc}=\frac{\dot Q_\mathrm{in,ORC}}
     {\left[h_\mathrm{w}(T_\mathrm{3d}+\Delta T_\mathrm{3D,2D})-h_\mathrm{w}(T_\mathrm{3d})\right]\Nw\,n_{t}},
  \label{eq:mdot}
\end{gather}
where $n_t$ is the number of hairpins per well and $h_\mathrm{w}(T)$ is the
specific enthalpy of water at $T$ and $p_\mathrm{w}$. Enthalpy differences are
used rather than a mean specific heat multiplied by the glide: over the
\SI{55}{\kelvin} charging rise the linearised form is \SI{0.1}{\percent}
inconsistent with the water enthalpies used elsewhere in the model. The
inconsistency is invisible to every energy balance, because both sides use the
same flow, but it generates a spurious exergy source of several kilowatts and
was detected by the exergy audit of Section~\ref{sec:exergy}.

The performance of the system is characterised by the indicators below, all
evaluated at cyclic steady state from the energies the field actually
transfers, $\Delta E_\mathrm{ch}$ and $\Delta E_\mathrm{dc}$, rather than from
the target values. The net electrical outputs are
\begin{equation}
  \dot W_\mathrm{el,out}^{*}=\frac{\Delta E_\mathrm{dc}}{t_\mathrm{dc}}\,
     \etaORC\,\eta_\mathrm{gen},
  \qquad
  \dot W_\mathrm{el,in}^{*}=\frac{\Delta E_\mathrm{ch}}{t_\mathrm{ch}\,
     \COP\,\eta_\mathrm{mot}},
  \label{eq:W-star}
\end{equation}
and, with $\dot W_\mathrm{f,ch}$ and $\dot W_\mathrm{f,dc}$ the pumping powers
of the whole field (Section~\ref{sec:pumping}),
\begin{align}
  \eta_{T}&=\frac{\dot W_\mathrm{el,out}^{*}\,t_\mathrm{dc}}{\Delta E_\mathrm{dc}}
           =\etaORC\,\eta_\mathrm{gen},
  &
  \eta_{T,\mathrm{eff}}&=\eta_{T}-\frac{\dot W_\mathrm{f,dc}\,t_\mathrm{dc}}
           {\Delta E_\mathrm{dc}},
  \label{eq:etaT}\\
  \mathrm{RTE}&=\frac{\left(\dot W_\mathrm{el,out}^{*}-\dot W_\mathrm{f,dc}\right)t_\mathrm{dc}}
                     {\left(\dot W_\mathrm{el,in}^{*}+\dot W_\mathrm{f,ch}\right)t_\mathrm{ch}},
  &
  \varepsilon_\mathrm{RTE}&=\mathrm{RTE}\left(\bar\varepsilon_\mathrm{ch}
           -\bar\varepsilon_\mathrm{dc}\right),
  \label{eq:rte}
\end{align}
where $\bar\varepsilon_\mathrm{ch}$ and $\bar\varepsilon_\mathrm{dc}$ are the
mean melted fractions of the store at the end of charging and of discharging
(Section~\ref{sec:state}). $\eta_{T}$ is a property of the power block alone;
$\eta_{T,\mathrm{eff}}$ charges the discharging pumping power against it;
RTE is the round-trip efficiency with both pumping parasitics; and
$\varepsilon_\mathrm{RTE}$ weights RTE by the fraction of the PCM inventory
that actually cycles, which penalises PCM that never changes phase. Weighting
by the cycled fraction rather than by the melted fraction at the end of
charging is necessary under the present formulation, because the melted
fraction saturates at unity once a segment is fully molten while further heat
continues to enter it as superheat. The thermal and net electrical energies
delivered per well and cycle, and the thermal energy density of the store, are
\begin{equation}
  \Delta E_\mathrm{therm}=\frac{\Delta E_\mathrm{dc}}{\Nw},\qquad
  \Delta E_\mathrm{elec}=\eta_{T,\mathrm{eff}}\,\Delta E_\mathrm{therm},\qquad
  \rho_{E}=\frac{\Delta E_\mathrm{therm}}{\Vwell}.
  \label{eq:dE}
\end{equation}

\paragraph{Exchanger conductance}
Every approach temperature in Table~\ref{tab:closure} improves an efficiency
at the expense of heat-exchanger area. To make that cost visible, the overall
conductance required by each of the four plant exchangers is evaluated from
its composite curves,
\begin{equation}
  UA=\int_{0}^{Q_\mathrm{tot}}\frac{\mathrm{d}Q}{\Delta T(Q)},
  \qquad
  \Delta T_\mathrm{eff}=\frac{Q_\mathrm{tot}}{UA},
  \label{eq:UA}
\end{equation}
where $\Delta T(Q)$ is the local hot-minus-cold temperature difference. For a
counterflow exchanger with two sensible streams of constant specific heat,
$\Delta T_\mathrm{eff}$ reduces to the logarithmic mean temperature
difference; with a phase change on one side it is its correct generalisation.

%-----------------------------------------------------------------------
\subsection{Borehole thermal regenerator}
\label{sec:BB}
%-----------------------------------------------------------------------

\subsubsection{Geometry and PCM cascade}
\label{sec:cascade}

Each well of depth $\Lwell$ and diameter $D_\mathrm{well}$ contains $n_t$
hairpins of finned tube, so that $2n_t$ tube legs run the full depth
(Fig.~\ref{fig:bb_geom}). The tubes have inner and outer radii $r_i$ and $r_e$
and carry $n_f$ longitudinal rectangular fins per leg, of radial length $L_f$
and thickness $t_f$. The developed length of one hairpin is
$\Ltube=2\Lwell$, and the coordinate $s\in[0,\Ltube]$ runs along it in the
direction of the charging flow. The remainder of the borehole is filled with
PCM, of volume
\begin{equation}
  \Vwell=\frac{\pi D_\mathrm{well}^{2}}{4}\Lwell
    -2n_t\left(\pi r_e^{2}+n_f t_f L_f\right)\Lwell .
  \label{eq:Vwell}
\end{equation}
The PCM belonging to one hairpin per unit developed length, and the radius of
the equivalent cylindrical cell owned by one tube leg, are
\begin{equation}
  \Aav=\frac{\Vwell}{n_t\,\Ltube},
  \qquad
  r_\mathrm{cell}=\frac{D_\mathrm{well}/2}{\sqrt{2n_t}} .
  \label{eq:Aav}
\end{equation}
$r_\mathrm{cell}$ is the radius at which the phase-change fronts of adjacent
legs meet, and $\delta_\mathrm{merge}=r_\mathrm{cell}-r_e$ is the largest
front thickness for which the annular representation of
Section~\ref{sec:network} is geometrically meaningful.

\begin{figure}[htbp]
  \centering
  \begin{subfigure}{0.42\linewidth}\centering
    \figorbox{\linewidth}{multilayer_BB_ch.png}
    \subcaption{Charging}
  \end{subfigure}\hfill
  \begin{subfigure}{0.42\linewidth}\centering
    \figorbox{\linewidth}{multilayer_BB_dc.png}
    \subcaption{Discharging}
  \end{subfigure}
  \caption{PCM layering in the borehole battery and the direction of the
  secondary-fluid flow in each half-cycle.}
  \label{fig:bb_geom}
\end{figure}

The PCM is divided along $s$ into $\Nlay$ layers of equal developed length,
each with its own melting temperature. During charging the water enters at
$s=0$, meets the layer of highest melting temperature first, and cools as it
proceeds; during discharging it enters at $s=\Ltube$, meets the layer of
lowest melting temperature first, and warms as it proceeds. The layers are
graded so that, for a linear water glide, each layer sits
$\Delta T_\mathrm{4C,M}$ below the charging water at its own leading face.
With a cascade parameter $\Delta T_\mathrm{casc}$ (equal to
$\Delta T_\mathrm{3C,2C}$ unless specified otherwise), the layer melting
temperatures and the coldest one are
\begin{gather}
  T_{m,k}=T_{m}^\mathrm{top}-(k-1)\frac{\Delta T_\mathrm{casc}}{\Nlay},
  \quad k=1,\ldots,\Nlay,
  \notag\\
  T_{m}^\mathrm{bot}\equiv T_{m,\Nlay}
  =T_{m}^\mathrm{top}-\Delta T_\mathrm{casc}\,\frac{\Nlay-1}{\Nlay}.
  \label{eq:cascade}
\end{gather}
The span of melting temperatures is one layer width short of
$\Delta T_\mathrm{casc}$, because the last layer is referenced to the water
entering it rather than leaving the well. Within each layer the local
superheat of the charging water therefore decays from
$\Delta T_\mathrm{4C,M}$ at the leading face to
$\Delta T_\mathrm{4C,M}-\Delta T_\mathrm{casc}/\Nlay$ at the trailing face.
The driving margin at the far end of the well, evaluated for a linear glide
against the melting temperatures, is positive only if
\begin{equation}
  \Delta T_\mathrm{casc}\frac{\Nlay-1}{\Nlay}
  >\Delta T_\mathrm{3C,2C}-\Delta T_\mathrm{4C,M}
  \label{eq:span-bound}
\end{equation}
during charging, and analogously with $\Delta T_\mathrm{3D,2D}$ and
$\Delta T_\mathrm{M,1D}$ during discharging. Equation~\eqref{eq:span-bound} is
a design guideline rather than a feasibility condition. It assumes that the
PCM sits at its melting temperature; under the enthalpy formulation of
Section~\ref{sec:state} the PCM superheats and subcools, the direction of heat
transfer is taken segment by segment from the sign of $T_j-\Tpcm$, and a
negative margin is handled consistently. Its consequence is that the far end
of the store is driven through a sensible branch, or briefly in reverse,
rather than through the phase change, which lowers the cycled fraction of the
inventory.

\subsubsection{Hypotheses}
\label{sec:hypotheses}

The borehole model rests on the following hypotheses.

\begin{description}[leftmargin=2.2em,itemsep=2pt,font=\normalfont\bfseries]
\item[H1] \emph{Lumped segment state.} The PCM belonging to a segment of
  length $\Delta s$ is represented by one enthalpy; its temperature is not
  resolved radially except through the closures of
  Section~\ref{sec:network}.
\item[H2] \emph{Quasi-steady conduction in the PCM.} The conduction
  resistance of the PCM is that of the steady profile for the instantaneous
  geometry. The approximation is controlled by the Stefan number
  $\mathrm{Ste}=c_{p}\Delta T/\hm$.
\item[H3] \emph{Concentric fronts.} The phase-change front about each leg is a
  concentric cylinder of radius $r_e+\delta$, unaffected by the fins and by
  the fronts of neighbouring legs, $\delta\le\delta_\mathrm{merge}$.
\item[H4] \emph{Conduction only.} Natural convection in the liquid PCM is
  neglected, consistent with a macro-encapsulated PCM, and there is no mixing
  between layers.
\item[H5] \emph{One front per segment.} A segment contains at most one
  solid--liquid interface.
\item[H6] \emph{Adiabatic store.} There is no heat exchange with the
  formation and no axial conduction in the PCM or the tube wall; segments
  interact only through the water.
\item[H7] \emph{Quasi-steady secondary fluid.} The water is an incompressible
  single-phase liquid whose heat capacity within the tube is negligible
  compared with the PCM; its temperature profile is established
  instantaneously at each time level. Viscous dissipation is neglected.
\item[H8] \emph{Developed-length cascade.} The layers are laid along the
  developed length of the hairpin, so the downward and upward legs carry
  different melting temperatures at the same depth.
\item[H9] \emph{Thermally independent legs.} No heat is exchanged between the
  two legs of a hairpin, or between hairpins.
\item[H10] \emph{Constant PCM volume.} The volume change on melting is
  neglected; the latent inventory is referred to the solid density $\rho_s$
  in both directions, so that the melted fraction is a mass fraction and the
  energy banked on charging equals that available on discharging.
\end{description}

\noindent
H3, H6 and H9 are the hypotheses most likely to affect the quantitative
results and are assessed in Section~\ref{sec:limitations}.

\subsubsection{Enthalpy state and constitutive relation}
\label{sec:state}

Let $E'$ denote the enthalpy of the PCM belonging to one segment, per unit
developed tube length, measured from a datum of fully solid material at that
segment's melting temperature $T_{m}$. The latent capacity and the two
sensible capacities per unit length are
\begin{equation}
  \Elat=\rho_s\hm\Aav,\qquad
  C_{s}=\rho_s c_{p,s}\Aav,\qquad
  C_{l}=\rho_s c_{p,l}\Aav .
  \label{eq:caps}
\end{equation}
The PCM temperature follows from $E'$ through the piecewise-linear
constitutive relation
\begin{equation}
  \Tpcm(E')=
  \begin{dcases}
    T_{m}+E'/C_{s}, & E'<0 \quad\text{(subcooled solid)},\\[2pt]
    T_{m},          & 0\le E'\le\Elat \quad\text{(two-phase)},\\[2pt]
    T_{m}+\left(E'-\Elat\right)/C_{l}, & E'>\Elat \quad\text{(superheated liquid)},
  \end{dcases}
  \label{eq:enth-curve}
\end{equation}
and the melted area and local melted fraction from
\begin{equation}
  \Amelt=\frac{\min\left(\max(E',0),\,\Elat\right)}{\rho_s\hm},
  \qquad
  \eloc=\frac{\Amelt}{\Aav}\in[0,1].
  \label{eq:amelt}
\end{equation}

Equation~\eqref{eq:enth-curve}, read as $E'(T)$, is multivalued at $T_m$,
which is why temperature cannot serve as the state variable; read as
$T(E')$, the direction used here, it is single-valued everywhere
(Fig.~\ref{fig:enthalpy-curve}). The saturation in Eq.~\eqref{eq:amelt} is
not a limiter added to suppress unphysical values but a consequence of the
definition: a segment cannot melt PCM it does not contain, and any further
energy is stored as superheat. The bound $\eloc\in[0,1]$ is therefore
satisfied by construction, and no heat delivered by the network is ever
discarded. A formulation that carries $\Amelt$ (or $\delta$) as its state
cannot represent the two sensible branches: once a segment has melted
completely it must either continue to melt non-existent material or discard
the heat it receives.

\begin{figure}[htbp]
\centering
\begin{tikzpicture}[>=Stealth, scale=0.9]
  \draw[->,thick] (0,0) -- (9.6,0) node[right] {$\Tpcm$};
  \draw[->,thick] (0,0) -- (0,5.8) node[above] {$E'$};
  \draw[dashed,gray] (4.9,0) -- (4.9,5.4);
  \node[below] at (4.9,0) {$T_{m}$};
  \draw[dashed,gray] (0,1.5) -- (4.9,1.5) node[pos=0,left] {$0$};
  \draw[dashed,gray] (0,3.5) -- (4.9,3.5) node[pos=0,left] {$\Elat$};
  \draw[very thick] (1.2,0.05) -- (4.9,1.5);
  \draw[very thick] (4.9,1.5) -- (4.9,3.5);
  \draw[very thick] (4.9,3.5) -- (8.9,5.05);
  \node[align=left,font=\small,anchor=west] at (0.85,2.65)
    {subcooled solid\\[1pt] slope $C_{s}$};
  \node[align=center,font=\small] at (7.0,5.35)
    {superheated liquid\\[1pt] slope $C_{l}$};
  \draw[<->,thick,black!60] (4.55,1.5) -- (4.55,3.5);
  \node[font=\small,align=right,anchor=east,black!70] at (4.4,1.95)
    {latent\\$\rho_s\hm\Aav$};
  \filldraw (4.9,2.6) circle (1.6pt);
  \draw[->,black!60] (5.9,1.75) -- (4.99,2.53);
  \node[font=\small,align=left,anchor=west,black!70] at (5.95,1.6)
    {$\eloc=E'/\Elat$};
\end{tikzpicture}
\caption{Constitutive relation of the segment state, Eq.~\eqref{eq:enth-curve}.
Read as $T(E')$ it is single-valued; position along the vertical segment is
the melted fraction.}
\label{fig:enthalpy-curve}
\end{figure}

\subsubsection{Segment energy balances and closure}
\label{sec:segment}

The developed length is divided into $n_s$ segments of length
$\Delta s=\Ltube/n_s$. For segment $j$, spanning $[s_j,s_{j+1}]$, two control
volumes share the tube wall (Fig.~\ref{fig:cv}): CV1, the PCM belonging to the
segment, which is closed and stationary; and CV2, the water within the tube
over the same interval, which enters at $T_j$ and leaves at $T_{j+1}$. Let
$q'_j$ denote the heat rate per unit length crossing from CV2 into CV1.

\begin{figure}[htbp]
\centering
\begin{tikzpicture}[>=Stealth, font=\small, scale=1.0]
  % PCM annulus (CV1)
  \fill[black!8] (0,1.1) rectangle (8,2.2);
  \fill[black!8] (0,-2.2) rectangle (8,-1.1);
  \draw[dashed] (0,1.1) rectangle (8,2.2);
  \draw[dashed] (0,-2.2) rectangle (8,-1.1);
  % melt shell
  \fill[orange!25] (0,1.1) rectangle (8,1.55);
  \fill[orange!25] (0,-1.55) rectangle (8,-1.1);
  % tube wall
  \fill[black!45] (0,0.85) rectangle (8,1.1);
  \fill[black!45] (0,-1.1) rectangle (8,-0.85);
  % water (CV2)
  \draw[thick,blue!60!black,dashed] (0,-0.8) rectangle (8,0.8);
  \draw[->,thick,blue!60!black] (-1.2,0) -- (-0.1,0) node[midway,above] {$\md,\,T_j$};
  \draw[->,thick,blue!60!black] (8.1,0) -- (9.3,0) node[midway,above] {$T_{j+1}$};
  \node[blue!60!black] at (4,0.25) {CV2: water};
  \node[blue!60!black] at (4,-0.3) {$\md c_p\,\mathrm{d}T/\mathrm{d}s=-q'$};
  \node at (4,1.9) {CV1: PCM, state $E'_j$, $\Tpcm(E'_j)$};
  \node at (4,-1.9) {CV1};
  \draw[->,thick,red!70!black] (2,0.55) -- (2,1.35) node[right,pos=0.7] {$q'_j$};
  \draw[->,thick,red!70!black] (6,-0.55) -- (6,-1.35) node[right,pos=0.7] {$q'_j$};
  % dimensions
  \draw[<->] (0,2.5) -- (8,2.5) node[midway,above] {$\Delta s$};
  \draw[<->] (8.4,0) -- (8.4,0.85);  \node[right] at (8.45,0.45) {$r_i$};
  \draw[<->] (9.2,0) -- (9.2,1.1);   \node[right] at (9.25,0.95) {$r_e$};
  \node[right,orange!70!black] at (8.1,1.4) {shell $\delta$};
  \draw[<->] (10.2,0) -- (10.2,2.2); \node[right] at (10.25,1.9) {$r_\mathrm{cell}$};
  \node[anchor=west] at (0,-2.55) {adiabatic at $r_\mathrm{cell}$ (H6, H9); no axial conduction (H6)};
\end{tikzpicture}
\caption{Control volumes of segment $j$ (axial section of one leg; fins not
shown). CV1 is the PCM belonging to the segment, CV2 the water over the same
interval. The shaded shell adjacent to the tube is the conduction path of
Eq.~\eqref{eq:shell-choice}.}
\label{fig:cv}
\end{figure}

\paragraph{PCM balance} CV1 does no work and, by H6, exchanges heat only
through the tube wall, so the first law reduces to
\begin{equation}
  \frac{\mathrm{d}E'_j}{\mathrm{d}t}=q'_j .
  \label{eq:pcm-balance}
\end{equation}

\paragraph{Water balance} For incompressible single-phase flow without
viscous dissipation, the energy equation per unit tube length is
\begin{equation}
  \rho_\mathrm{w}A_\mathrm{flow}c_{p}\frac{\partial T}{\partial t}
  +\md c_{p}\frac{\partial T}{\partial s}=-q'(s),
  \label{eq:fluid-pde}
\end{equation}
with $A_\mathrm{flow}=\pi r_i^{2}$. Hypothesis H7 drops the storage term. The
approximation rests on the separation between the transit time of the water
through the hairpin, $\Ltube\rho_\mathrm{w}A_\mathrm{flow}/\md$, and the time
over which the PCM state evolves, $t_\mathrm{ch}$ or $t_\mathrm{dc}$, and on
the ratio of the energy held in the water to the latent inventory,
$\rho_\mathrm{w}A_\mathrm{flow}c_{p}\Delta T/\Elat$; both are quantified for
the reference case in Section~\ref{sec:conditions}. Its cost is that the
start-up transient of about one transit time at the beginning of each
half-cycle is not resolved.

The local heat flux is closed on the resistance network of
Section~\ref{sec:network}, referred to the inner tube area,
\begin{equation}
  q'(s)=2\pi r_i\,\Ui\left[T(s)-\Tpcm\right],
  \label{eq:local-flux}
\end{equation}
where CV1 presents the single temperature $\Tpcm(E'_j)$ over the whole
segment (H1). The perimeter is the inner perimeter $2\pi r_i$ because $\Ui$ is
referred to the inner area and already contains the finned outer surface.

\paragraph{Integration across the segment} With $\Tpcm$ and $\Ui$ held
constant over the segment, the steady form of Eq.~\eqref{eq:fluid-pde} with
Eq.~\eqref{eq:local-flux} is a linear ordinary differential equation in $s$
whose solution is the single-stream exchanger result
\begin{equation}
  T_{j+1}=\Tpcm+\left(T_j-\Tpcm\right)e^{-\NTU_j},
  \qquad
  \NTU_j=\frac{2\pi r_i\,\Ui\,\Delta s}{\md c_{p}} .
  \label{eq:ntu}
\end{equation}
The water is a stream of finite capacity rate $\md c_p$ exchanging heat with
an isothermal reservoir; the effectiveness is $1-e^{-\NTU}$.

\paragraph{Closure} The heat given up by the water over the segment is the
enthalpy it loses, $\md c_p(T_j-T_{j+1})=q'_j\Delta s$, so that
\begin{equation}
  q'_j=K_j\left(T_j-\Tpcm\right),
  \qquad
  K_j=\frac{\md c_{p}\left(1-e^{-\NTU_j}\right)}{\Delta s},
  \label{eq:K}
\end{equation}
and substituting into Eq.~\eqref{eq:pcm-balance} gives the segment equation
\begin{equation}
  \frac{\mathrm{d}E'_j}{\mathrm{d}t}=K_j\left[T_j-\Tpcm(E'_j)\right],
  \label{eq:ode}
\end{equation}
with one unknown per segment; $T_j$ is handed down by the upstream segment and
$K_j$ depends on the state through $\Ui$. The water temperature is then
advanced by
\begin{equation}
  T_{j+1}=T_j-\frac{q'_j\,\Delta s}{\md c_{p}} .
  \label{eq:T-march}
\end{equation}

The closure of Eq.~\eqref{eq:K} is not the product
$2\pi r_i\Ui$ that results from evaluating Eq.~\eqref{eq:local-flux} at the
segment inlet. The ratio of the two is
\begin{equation}
  \frac{K_j}{2\pi r_i\Ui}=\frac{1-e^{-\NTU_j}}{\NTU_j}\le 1 .
  \label{eq:K-ratio}
\end{equation}
As $\NTU\to0$ the ratio tends to unity; as $\NTU\to\infty$, $K\to\md
c_p/\Delta s$, a limit set by the capacity rate of the water rather than by the
wall. The product form has no such limit and permits $T_{j+1}<\Tpcm$, i.e.\
heat transfer from cold to hot, at a rate that increases without bound as the
wall conductance improves. The discrepancy is largest where the phase-change
shell is thin and the exchanger most effective.

\subsubsection{Thermal resistance network}
\label{sec:network}

The overall heat transfer coefficient referred to the inner tube area is the
series combination of internal convection, internal fouling, wall conduction
and the PCM-side resistance of the finned outer surface,
\begin{equation}
  \frac{1}{\Ui}=\frac{1}{h_i}+R''_{f,i}+\frac{r_i}{k_\mathrm{wall}}\ln\frac{r_e}{r_i}
  +\frac{2\pi r_i}{P_T\,\etao\,h_e},
  \qquad
  h_e=\frac{k}{r_e\ln\left(1+\delta/r_e\right)},
  \label{eq:Ui}
\end{equation}
where $h_e$ is the conductance of a cylindrical PCM shell of thickness $\delta$
and conductivity $k$, referred to the outer tube surface, and the factor
$2\pi r_i/(P_T\etao)$ transfers it to the inner area.

\paragraph{Internal convection} The coefficient $h_i$ is evaluated locally at
the bulk water temperature $T_j$ of each segment and at each time level. For
$\mathrm{Re}=4\md/(\pi D_i\mu)\ge2300$ the Gnielinski correlation
\cite{Gnielinski1976} is used with the Petukhov friction factor,
\begin{equation}
  \mathrm{Nu}=\frac{(f/8)(\mathrm{Re}-1000)\mathrm{Pr}}
    {1+12.7\sqrt{f/8}\left(\mathrm{Pr}^{2/3}-1\right)},
  \qquad
  f=\left(0.79\ln\mathrm{Re}-1.64\right)^{-2},
  \label{eq:gnielinski}
\end{equation}
and $\mathrm{Nu}=3.66$ otherwise, with $h_i=\mathrm{Nu}\,k_\mathrm{w}/D_i$, where $k_\mathrm{w}$ is the conductivity of water. Properties are evaluated at $T_j$ and $p_\mathrm{w}$.

\paragraph{Fins} With $n_f$ longitudinal fins per leg, the finned perimeter,
the fin efficiency and the overall surface efficiency are
\begin{gather}
  P_T=2\pi r_e+2n_fL_f,
  \qquad
  \etao=1-\frac{n_fA_f}{A_t}\left(1-\etaf\right),
  \label{eq:fins}\\
  \etaf=\frac{\tanh(mL_c)}{mL_c},
  \qquad m=\sqrt{\frac{2h_e}{k_\mathrm{fin}t_f}},
  \qquad L_c=L_f+\frac{t_f}{2},
  \label{eq:fin-eff}
\end{gather}
with $A_f=(2L_f+t_f)\Delta s$ the surface of one fin and $A_t=P_T\Delta s$
the total outer surface of the segment \cite{Incropera}. The fin conductivity
is taken equal to that of the tube wall, $k_\mathrm{fin}=k_\mathrm{wall}$.

\paragraph{Shell thickness from the melted area}
The fin metal occupies part of the annulus about the tube, so the melted area
and the shell thickness are related by
\begin{equation}
  A=\pi\left[(r_e+\delta)^{2}-r_e^{2}\right]-n_ft_f\min(\delta,L_f),
  \label{eq:area-fins}
\end{equation}
which is inverted in closed form on each branch:
\begin{align}
  \delta\le L_f:&\quad
    \pi\delta^{2}+\left(2\pi r_e-n_ft_f\right)\delta-A=0,
  \label{eq:delta-in}\\
  \delta>L_f:&\quad
    \pi\delta^{2}+2\pi r_e\delta-\left(A+n_ft_fL_f\right)=0 .
  \label{eq:delta-out}
\end{align}

\paragraph{Conduction path: which shell the heat crosses}
The PCM-side term in Eq.~\eqref{eq:Ui} is the resistance of a shell growing
outward from the tube wall, so the model must specify the material of which
that shell consists. Because the tube is the driven boundary in both
half-cycles, the phase-change front always originates at the wall and moves
outward. During melting the shell is liquid and its area is the melted area;
during freezing the shell is the newly frozen solid and its area is the solid
area, the remaining liquid being displaced beyond the front and out of the
conduction path. For a segment in the two-phase branch,
\begin{equation}
  \left(A,\,k\right)=
  \begin{cases}
    \left(\Amelt,\;k_{l}\right), & T_j>\Tpcm\quad\text{(melting)},\\[3pt]
    \left(\Aav-\Amelt,\;k_{s}\right), & T_j<\Tpcm\quad\text{(freezing)},
  \end{cases}
  \label{eq:shell-choice}
\end{equation}
and $\delta$ follows from Eqs.~\eqref{eq:delta-in}--\eqref{eq:delta-out}. The
direction is taken from the sign of the driving temperature difference, which
is known before the conductance is formed. Equation~\eqref{eq:shell-choice}
gives the correct limits without special treatment: a fully solid segment
beginning to melt and a fully liquid segment beginning to freeze both have
$\delta=0$. Using the melted area in both directions, as in a formulation
written for melting and not revisited, reverses the freezing half-cycle in
time: a segment frozen solid is assigned zero resistance where the physical
resistance is greatest.

\paragraph{Single-phase segments: a mean-to-surface resistance}
A segment in either sensible branch contains no front, and
Eq.~\eqref{eq:shell-choice} does not apply. In that case $\Tpcm$ is the
volume-mean temperature of the cell, because the state is the cell's total
enthalpy, and the resistance consistent with a mean temperature is the
mean-to-surface resistance. Consider the cell as an annulus
$r_e\le r\le r_\mathrm{cell}$, adiabatic at $r_\mathrm{cell}$ (a symmetry
surface with the neighbouring cell), storing sensible heat at a uniform
volumetric rate $\dot s$. Under quasi-steady conditions (H2) the heat crossing
radius $r$ is that stored beyond it,
$-k\,2\pi r\,\mathrm{d}T/\mathrm{d}r=\dot s\pi(r_\mathrm{cell}^{2}-r^{2})$,
which integrates to
\begin{equation}
  T(r_e)-T(r)=\frac{\dot s}{2k}
  \left[r_\mathrm{cell}^{2}\ln\frac{r}{r_e}-\frac{r^{2}-r_e^{2}}{2}\right].
  \label{eq:bulk-profile}
\end{equation}
Averaging over the cell volume and dividing by
$Q'=\dot s\pi(r_\mathrm{cell}^{2}-r_e^{2})$ gives
$R'_\mathrm{bulk}=S/(2\pi k)$, with the geometric shape factor
\begin{equation}
  S=\frac{4\beta^{4}\ln\beta-3\beta^{4}+4\beta^{2}-1}
         {4\left(\beta^{2}-1\right)^{2}},
  \qquad \beta=\frac{r_\mathrm{cell}}{r_e}.
  \label{eq:shape-factor}
\end{equation}
The conductivity cancels, so the same $S$ serves both phases. For a thin
annulus of thickness $L=r_\mathrm{cell}-r_e\ll r_e$, $S\to L/(3r_e)$ and
$R'_\mathrm{bulk}\to L/(3kA)$ with $A=2\pi r_e$ per unit length, which is the
mean-to-surface resistance of a plane slab with uniform generation. The
resistance enters Eq.~\eqref{eq:Ui} as an equivalent thickness
$\delta_\mathrm{eq}=r_e\left(e^{S}-1\right)$, with $k=k_s$ for subcooled
solid and $k=k_l$ for superheated liquid, so that the fin relations apply
unchanged.

The PCM-side resistance is consequently discontinuous at the branch
boundaries. As $\eloc\to1^{-}$ during melting the front approaches
$r_\mathrm{cell}$ and $R'\to\ln\beta/(2\pi k_l)$, whereas once the segment is
superheated $R'=S/(2\pi k_l)$. The discontinuity reflects the change in the
reference temperature, from the front to the volume mean, rather than a
physical event, and the heat flux remains bounded because $K$ is limited by the
capacity rate, Eq.~\eqref{eq:K}.

\paragraph{Limitation of a single front} After the first half-cycle from a
uniformly solid store, the phase distribution within a cell is in general
three-region: at cyclic steady state a charge begins with residual liquid in
the outer part of the cell, so that melting produces liquid at the tube, solid
in the middle and liquid beyond it. Equation~\eqref{eq:shell-choice} is the
better of the two single-front approximations, not a resolution of this
geometry; this is the content of H5.

%-----------------------------------------------------------------------
\subsection{Charging period}
\label{sec:BB_ch}
%-----------------------------------------------------------------------

During charging the water enters every hairpin at $s=0$ with temperature
$T_\mathrm{4c}$, Eq.~\eqref{eq:inlets}, and mass flow rate $\md_\mathrm{ch}$,
Eq.~\eqref{eq:mdot}, and the segment equations
\eqref{eq:ode}--\eqref{eq:T-march} are marched in the direction of increasing
$s$ against the cascade of Eq.~\eqref{eq:cascade}. Segments in which
$T_j>\Tpcm$ melt through a liquid shell of conductivity $k_l$; segments that
have completed melting superheat, their driving difference $T_j-\Tpcm$
decreases, and the heat they no longer absorb is carried downstream. The
energy absorbed by the field over the charging period and the melted fraction
of the store at its end are
\begin{gather}
  \Delta E_\mathrm{ch}=\Nw\,n_t\int_{0}^{t_\mathrm{ch}}\sum_{j=1}^{n_s}
     q'_j(t)\,\Delta s\,\mathrm{d}t
   =\Nw\,n_t\,\Delta s\sum_{j=1}^{n_s}\left[E'_j(t_\mathrm{ch})-E'_j(0)\right],
  \label{eq:Ech}\\
  \bar\varepsilon_\mathrm{ch}=\frac1{n_s}\sum_{j=1}^{n_s}
     \varepsilon_{\mathrm{loc},j}(t_\mathrm{ch}),
  \label{eq:eps-ch}
\end{gather}
where the second equality holds exactly under the integrator of
Section~\ref{sec:integration}. The water leaves the field at the
time-dependent temperature $T_{n_s+1}(t)$; its flow-averaged value,
\begin{equation}
  \overline{T}_\mathrm{2c}=\frac{1}{t_\mathrm{ch}}\int_{0}^{t_\mathrm{ch}}
     T_{n_s+1}(t)\,\mathrm{d}t,
  \label{eq:Tout-mean}
\end{equation}
is the temperature at which the water returns to the HTHP condenser and is
reported alongside the nominal value $T_\mathrm{2c}$ of
Table~\ref{tab:closure}. At constant $\md$ and $c_p$ the flow average reduces
to the time average; on a logarithmic time grid it must be weighted by the
step length.

%-----------------------------------------------------------------------
\subsection{Discharging period}
\label{sec:BB_dc}
%-----------------------------------------------------------------------

During discharging the flow reverses. The water enters at $s=\Ltube$ with
temperature $T_\mathrm{3d}$, Eq.~\eqref{eq:inlets}, and mass flow rate
$\md_\mathrm{dc}$, and meets the layer of lowest melting temperature first.
The governing equations are those of the charging period, with two
differences. First, in segments where $T_j<\Tpcm$ the conduction path is the
frozen shell with conductivity $k_s$, Eq.~\eqref{eq:shell-choice}. Second, the
march proceeds in the direction of decreasing $s$, so both the melting
temperatures and the initial state, which is the state at the end of the
preceding charge, are mirrored with the march index; since $E'$ is measured
from a datum at each segment's own $T_m$, mirroring the cascade without
mirroring the state would displace the datum by up to the full cascade span.
The energy released by the field and the residual melted fraction are
\begin{equation}
  \Delta E_\mathrm{dc}=-\Nw\,n_t\,\Delta s\sum_{j=1}^{n_s}
     \left[E'_j(t_\mathrm{dc})-E'_j(0)\right],
  \qquad
  \bar\varepsilon_\mathrm{dc}=\frac1{n_s}\sum_{j=1}^{n_s}
     \varepsilon_{\mathrm{loc},j}(t_\mathrm{dc}),
  \label{eq:Edc}
\end{equation}
and the flow-averaged outlet $\overline{T}_\mathrm{2d}$, defined as in
Eq.~\eqref{eq:Tout-mean}, is the temperature at which the water reaches the ORC
evaporator.

Idle periods between the half-cycles, if any, leave the state unchanged under
H6, since the store is adiabatic and carries no axial conduction.

%-----------------------------------------------------------------------
\subsection{Pressure drop and pumping power}
\label{sec:pumping}
%-----------------------------------------------------------------------

The pressure drop over one hairpin is the sum of distributed friction over the
developed length and local losses at the entrance, the return bends and the
exit,
\begin{equation}
  \Delta p_f=\left(f\,\frac{\Ltube}{D_i}+\sum K_\mathrm{loc}\right)
            \frac{\rho_\mathrm{w}\bar u^{2}}{2},
  \qquad
  \bar u=\frac{\md}{\rho_\mathrm{w}\pi r_i^{2}},
  \label{eq:dp}
\end{equation}
with $\sum K_\mathrm{loc}=0.5+2(2.2)+1.0$ \chk{one 180$^\circ$ return per
hairpin, or two?} and the Darcy friction factor of a smooth tube:
$f=64/\mathrm{Re}$ for $\mathrm{Re}<2000$; $f=0.3164\,\mathrm{Re}^{-1/4}$
(Blasius) for $\mathrm{Re}\le10^{5}$; and
$f=0.0032+0.221\,\mathrm{Re}^{-0.237}$ for $10^{5}<\mathrm{Re}\le10^{6}$.
\chk{Ref.~\cite{KoIHTC2026} states Haaland \cite{Haaland1983}; the code uses the
smooth-tube forms above. Choose one and make paper and code agree.} Water
properties are evaluated at the loop pressure and the mean of the inlet and
the flow-averaged outlet temperatures. The hydraulic pumping power of the
field in each half-cycle is
\begin{equation}
  \dot W_{f}=\Nw\,n_t\,\frac{\md}{\rho_\mathrm{w}}\,\Delta p_f ,
  \label{eq:Wf}
\end{equation}
\chk{no pump efficiency is applied; add $\eta_\mathrm{pump}$ or state it.}
evaluated with $\md_\mathrm{ch}$ and $\md_\mathrm{dc}$ respectively and
entering Eqs.~\eqref{eq:etaT}--\eqref{eq:rte}. Hydrostatic contributions
cancel over a closed hairpin.

